\section*{第 \ref{ch:b1:complexation} 章——配合作用}

\begin{solution}{exo:b1:complexation:1}
$[\ce{Cu(NH3)4}]^{2+}$：\ce{Cu^2+}，四个 \ce{NH3} 以 N 配位。
$[\ce{Fe(CN)6}]^{4-}$：\ce{Fe^2+}，六个 \ce{CN-} 以 C 配位；
$[\ce{Fe(CN)6}]^{3-}$：与上同，中心为 \ce{Fe^3+}。$[\ce{Ag(NH3)2}]^{+}$：
\ce{Ag+}，两个 \ce{NH3} 以 N 配位。$[\ce{FeSCN}]^{2+}$：\ce{Fe^3+}，一个
\ce{SCN-}，尽管化学式这样书写，仍以 N 配位。
$[\ce{Zn(OH)4}]^{2-}$：\ce{Zn^2+}，四个 \ce{OH-} 以 O 配位。
$[\ce{CaY}]^{2-}$：\ce{Ca^2+}，一个 EDTA 离子以两个 N 和四个 O 配位。
\end{solution}

\begin{solution}{exo:b1:complexation:2}
$\log K_i = \log\beta_i - \log\beta_{i-1}$：依次为 4.10、3.41、2.82 和 2.03，
$\mathrm{p}K_{d,i} = \log K_i$ 取相同的值。最后一个对应的是
$[\ce{Cu(NH3)4}]^{2+} \rightleftharpoons [\ce{Cu(NH3)3}]^{2+} + \ce{NH3}$.
\end{solution}

\begin{solution}{exo:b1:complexation:3}
优势区域：$[\ce{Cu(NH3)4}]^{2+}$ 位于 pNH3 2.03 以下，其后 $n = 3$ 直至 2.82，
$n = 2$ 直至 3.41，$n = 1$ 直至 4.10，\ce{Cu^2+} 居其上。pNH3 = 0 时为
$[\ce{Cu(NH3)4}]^{2+}$；2.5 时为 $[\ce{Cu(NH3)3}]^{2+}$；3.5 时为
$[\ce{CuNH3}]^{2+}$；5.0 时为 \ce{Cu^2+}。
\end{solution}

\begin{solution}{exo:b1:complexation:4}
$[\ce{FeSCN^2+}]/[\ce{Fe^3+}] = \beta_1[\ce{SCN-}] = 10^{2.96} \times 0.10 =
91$：配合分数为 $91/92 = \qty{98.9}{\percent}$。
\end{solution}

\begin{solution}{exo:b1:complexation:5}
每个 $[\mathrm{ML}_k] = \beta_k[\mathrm{M}][\mathrm{L}]^k$；两边除以
对 $k$ 求和（即金属总量），可消去 $[\mathrm{M}]$。若有三种
物种，$f_i = 1/(1 + [\mathrm{ML}_{i-1}]/[\mathrm{ML}_i] +
[\mathrm{ML}_{i+1}]/[\mathrm{ML}_i]) = 1/(1 + 1/(K_i x) + K_{i+1}x)$ 其中
$x = [\mathrm{L}]$。分母最小时，其导数
$-1/(K_i x^2) + K_{i+1}$ 为零，即 $x^2 = 1/(K_iK_{i+1})$，此时
两邻近项相等，且 $\mathrm{pL} = \frac12(\log K_i + \log K_{i+1})$。
对 $[\ce{Cu(NH3)2}]^{2+}$：$\frac12(3.41 + 2.82) = 3.12$。
\end{solution}

\begin{solution}{exo:b1:complexation:6}
$[\ce{FeSCN}]^{2+} + \ce{F-} \rightleftharpoons [\ce{FeF}]^{2+} + \ce{SCN-}$,
$K = 10^{6.85 - 2.96} = 10^{3.89}$。比值 $= K[\ce{F-}]/[\ce{SCN-}] =
10^{3.89} \times 0.010/0.10 = 780$：红色\omterm{def:b1:complexation:complex}{配合物}几乎全部
被取代，颜色褪去。
\end{solution}

\begin{solution}{exo:b1:complexation:7}
$[\ce{MgY}]^{2-} + \ce{Ca^2+} \rightleftharpoons [\ce{CaY}]^{2-} + \ce{Mg^2+}$,
$K = 10^{12.69 - 10.90} = 10^{1.79} = 62$。若两者等量，
$x^2/(1 - x)^2 = 62$，$x/(1 - x) = 7.9$，$x = 0.89$：\qty{89}{\percent} 的
镁被释放出来。
\end{solution}

\begin{solution}{exo:b1:complexation:8}
\ce{Ag2O(s) + H2O + 4NH3 <=> 2[Ag(NH3)2]+ + 2OH-}, $K = \beta_2^2 \times
10^{-15.43} = 10^{14.44 - 15.43} = 10^{-0.99}$。常数并不
很小：只要游离氨达到零点几摩尔每升，
首批几滴生成的棕色氧化物（氨在此作
碱）溶解生成银氨\omterm{def:b1:complexation:complex}{配合物}。
\end{solution}

\begin{solution}{exo:b1:complexation:9}
pH 10 时：$\alpha = 1 + 10^{1.24} + 10^{-1.96} = 18.4$，$\log\beta' = 12.69 -
1.26 = 11.43$。pH 7 时：$\alpha = 1 + 10^{4.24} + 10^{4.04} + 10^{0.19} +
\cdots = 10^{4.45}$，$\log\beta' = 8.24$。在 pH 7 时，条件常数
低于 $10^{10}$：钙与 EDTA 的反应不够完全，
不足以用于滴定，因为大部分 EDTA 已质子化。
\end{solution}

\begin{solution}{exo:b1:complexation:10}
1.~\ce{CuO(s) + H2O + 4NH3 <=> [Cu(NH3)4]^2+ + 2OH-}, $K = \beta_4 \times
10^{-20.64} = 10^{12.36 - 20.64} = 10^{-8.28}$.
2.~$\mathrm{pH} = 7 + \frac12(9.25 + \log 1.0) = 11.63$, $[\ce{OH-}] =
10^{-2.37}$，且 $[\ce{Cu(NH3)4^2+}] = 10^{-8.28}/10^{-4.75} =
\qty{3.0e-4}{mol/L}$：溶解量很少。
3.~$\mathrm{pH} = 9.25$，$[\ce{OH-}] = 10^{-4.75}$，由该式得
$10^{-8.28 + 9.50} = 10^{1.22}$，远超氨所能提供的量：
氧化物会不断溶解，直到氨用尽为止，而不是达到平衡。
铵离子消耗溶解过程中释放的氢氧根离子，这
正是铵盐有助于氨溶解铜化合物的原因。
\end{solution}

\begin{solution}{exo:b1:complexation:11}
1.~该反应为第二步生成反应减去第一步：$K = K_2/K_1 =
10^{3.90 - 3.32} = 10^{0.58} = 3.8 > 1$，故 $[\ce{AgNH3}]^{+}$ 会与
自身发生反应，即当其为主要物种时。
2.~由\cref{exo:b1:complexation:5}可知，最大分数位于
$\mathrm{pNH_3} = \frac12(3.32 + 3.90) = 3.61$ 处，此时 $\beta_1 x = 10^{-0.29}
= 0.51$，$\beta_2x^2 = 1.0$：$f = 0.51/(1 + 0.51 + 1.0) = 0.20$。从未
超过 \qty{20}{\percent} 的银以 $[\ce{AgNH3}]^{+}$ 形式存在。
\end{solution}

\begin{solution}{exo:b1:complexation:12}
1.~\ce{CaCO3(s) + Y^4- <=> [CaY]^2- + CO3^2-}, $K = \beta K_s = 10^{12.69 -
8.30} = 10^{4.39}$.
2.~$x^2/(0.10 - x) = 10^{4.39}$ 给出 $x = \qty{0.10}{mol/L}$，误差在
$4 \times 10^{-7}$ 之内：EDTA 全部用完，
相当于每升 \qty{10.0}{g} 碳酸钙。
3.~在酸性条件下，EDTA 被质子化（$\mathrm{p}K_{a4} = 11.24$），其
条件常数下降；碳酸根也被质子化，这有所帮助，
但固定钙的是\omterm{def:b1:complexation:complex}{配合物}。保持较高的 pH，使
EDTA 保持为 \ce{Y^4-}。
\end{solution}

\begin{solution}{pb:b1:complexation:1}
\textbf{1.} \ce{Ag+ + NH3 <=> [AgNH3]+}, $\beta_1 = [\ce{AgNH3+}]/([\ce{Ag+}][\ce{NH3}])$;
\ce{Ag+ + 2NH3 <=> [Ag(NH3)2]+}, $\beta_2 = [\ce{Ag(NH3)2+}]/([\ce{Ag+}][\ce{NH3}]^2)$.
\textbf{2.} $\log K_1 = 3.32$, $\log K_2 = 7.22 - 3.32 = 3.90$.
\textbf{3.} $[\ce{AgNH3}]^{+}$ 将在 pNH3 3.90 与 3.32 之间占优势，
该区间是倒置的，因此没有优势区域。
\textbf{4.} \ce{2[AgNH3]+ <=> Ag+ + [Ag(NH3)2]+}, $K = K_2/K_1 = 10^{0.58} =
3.8$.
\textbf{5.} 以下为 $[\ce{Ag(NH3)2}]^{+}$，以上为 \ce{Ag+}，分界点为
$\frac12\log\beta_2 = 3.61$.
\textbf{6.} $\beta_2[\ce{NH3}]^2 = 10^{7.22 - 2.00} = 10^{5.22} = \num{1.7e5}$.
\textbf{7.} $\mathrm{p}K_d = \log\beta_2 = 7.22$.
\textbf{8.} \ce{AgCl(s) + 2NH3 <=> [Ag(NH3)2]+ + Cl-}, $K = \beta_2K_s =
10^{7.22 - 9.75} = 10^{-2.53}$.
\textbf{9.} $s = \sqrt{K}\,[\ce{NH3}] = 10^{-1.265} \times 1.0 =
\qty{0.054}{mol/L}$.
\textbf{10.} 在水中 $s = 10^{-4.875} = \qty{1.3e-5}{mol/L}$：约
多出 4000 倍。
\textbf{11.} $0.0543 \times 143.4 = \qty{7.8}{g}$.
\textbf{12.} $[\ce{Ag+}] = K_s/[\ce{Cl-}] = 10^{-9.75}/0.054 =
\qty{3.3e-9}{mol/L} \ll s$.
\textbf{13.} $s = \sqrt{K}(1.0 - 2s)$, $s = 0.0543/(1 + 2 \times 0.0543) =
\qty{0.049}{mol/L}$.
\textbf{14.} \qty{1.0}{mol/L} 氨水的 pH 为 11.6，大于 $9.25 + 1$：
只有 \qty{0.4}{\percent} 的氨以 \ce{NH4+} 形式存在。
\textbf{15.} $K = 10^{7.22 - 12.27} = 10^{-5.05}$.
\textbf{16.} $s = 10^{-2.525} = \qty{3.0e-3}{mol/L}$, \qty{0.56}{g/L}.
\textbf{17.} $K = 10^{-8.85}$, $s = 10^{-4.425} = \qty{3.8e-5}{mol/L}$,
\qty{8.8}{mg/L}.
\textbf{18.} 分别用稀氨水处理各沉淀：只有氯化物
溶解。浓氨水溶解溴化物，速度较慢，且
仅部分溶解；碘化物不溶解。
\textbf{19.} $s = 1.0/187.8 = \qty{5.3e-3}{mol/L}$; $[\ce{NH3}] = s/\sqrt{K}
= \num{5.3e-3}/10^{-2.525} = \qty{1.8}{mol/L}$.
\textbf{20.} $s = \qty{4.3e-3}{mol/L}$，且 $[\ce{NH3}] = \num{4.3e-3}/10^{-4.425}
= \qty{113}{mol/L}$，是纯水中水浓度的
两倍（\qty{55}{mol/L}）：不可能。氨不能溶解碘化银。
\textbf{21.} \ce{[Ag(NH3)2]+ + Cl- + 2H3O+ -> AgCl(s) + 2NH4+ + 2H2O},
$K = 10^{2.53} \times (10^{9.25})^2 = 10^{21.03}$：白色沉淀
重新出现，证实为氯化银。
\textbf{22.} $s = 1.0/143.4 = \qty{6.97e-3}{mol/L}$; $[\ce{NH3}] =
s/\sqrt{K} = \num{6.97e-3}/10^{-1.265} = \qty{0.128}{mol/L}$.
\textbf{23.} $2s = \qty{0.0139}{mol/L}$.
\textbf{24.} $[\ce{Ag+}] = 10^{-9.75}/\num{6.97e-3} = \qty{2.6e-8}{mol/L}$,
与 $s$ 相比可忽略。
\textbf{25.} $0.128 + 0.014 = $ \textbf{\qty{0.142}{mol/L} 的氨}：
\qty{0.142}{mol} 的氨在一升溶液中恰好溶解 \qty{1.0}{g} 的
氯化银。
\end{solution}
